\begin{document}
\begin{center}
{\Large\textbf{Конус, цилиндр и шар: постоянные отношения}}\\[2mm]
{\small ЕГЭ, стереометрия, задание №3 \quad | \quad Фамилия\ \answerline{6cm}}
\end{center}
\ifteacher
\begin{center}\fbox{\textbf{Учительская версия:} ответы и краткие опоры выделены курсивом.}\end{center}
\fi

\textbf{Цель.} Строить осевое сечение и выводить отношение из формул, а не
запоминать отдельный ответ. Рисунок является схемой: измерять на нём ничего не нужно.

\section*{1. Общие основание и высота}
\stereotask{1}{Цилиндр и конус имеют общие основание и высоту. Объём цилиндра
равен \(162\). Найдите объём конуса.}{\csname figStereoV21\endcsname}
{Объём конуса составляет треть объёма цилиндра: \(54\).}

\stereotask{2}{Цилиндр и конус имеют общие основание и высоту. Высота цилиндра
равна радиусу основания. Площадь боковой поверхности цилиндра равна
\(27\sqrt2\). Найдите площадь боковой поверхности конуса.}
{\csname figStereoV22\endcsname}
{При \(h=r\) образующая \(r\sqrt2\); площадь боковой поверхности конуса равна \(27\).}

\section*{2. Конус и шар}
\stereotask{3}{Конус вписан в шар. Радиус основания конуса равен радиусу шара.
Объём шара равен \(188\). Найдите объём конуса.}{\csname figStereoV23\endcsname}
{В этом случае объём шара в четыре раза больше объёма конуса. Ответ: \(47\).}

\stereotask{4}{Конус вписан в шар. Радиус основания конуса равен радиусу шара.
Объём конуса равен \(19\). Найдите объём шара.}{\csname figStereoV24\endcsname}
{Объём шара в четыре раза больше. Ответ: \(76\).}

\section*{3. Шар в цилиндре}
\stereotask{5}{Шар объёмом \(64\) вписан в цилиндр. Найдите объём цилиндра.}
{\csname figStereoV29\endcsname}
{У вписанного шара высота цилиндра равна \(2R\), поэтому его объём равен
\(3\cdot64/2=96\).}

\stereotask{6}{Шар вписан в цилиндр. Площадь поверхности шара равна \(74\).
Найдите площадь полной поверхности цилиндра.}{\csname figStereoV30\endcsname}
{Площадь полной поверхности цилиндра составляет \(3/2\) площади шара. Ответ: \(111\).}

\stereotask{7}{Шар вписан в цилиндр. Площадь поверхности шара равна \(26\).
Найдите площадь полной поверхности цилиндра.}{\csname figStereoV32\endcsname}
{Площадь полной поверхности цилиндра составляет \(3/2\) площади шара. Ответ: \(39\).}

\ifteacher
\section*{Короткая проверка урока}
Ученик готов двигаться дальше, если отличает радиус шара от диаметра и может
вывести отношения \(1:3\), \(1:4\) и \(3:2\) через осевое сечение.
\fi
\end{document}
